\section{Introduction}
\label{sec:introduction}

Federated graph learning lets institutions train graph neural networks without pooling their subgraphs~\citep{kipf2016semi,velickovic2017graph,mcmahan2017communication,kairouz2021advances}. In credit and risk applications, however, collaboration must address both group disparities and malicious participants. Several fairness-aware aggregation rules weight clients using locally reported losses or fairness statistics~\citep{ezzeldin2023fairfed,li2020fair,f2gnn2023,khan2026fedgraphfair,bendoukha2025popets,dang2026fedfairgnn}. Those statistics may be useful under honest participation, but the server cannot usually verify them without access to private client data. Consider a consortium of lenders that trains a credit-risk model on its members' customer graphs. If one member can raise its own aggregation weight by reporting a small disparity and a low loss, it can steer the shared model while appearing to be the most responsible participant, and no other member can detect this from the reports alone.

This paper asks a concrete operational question: what happens when an aggregation rule treats a self-reported scalar as evidence of client quality? We first separate that interface dependency from training dynamics. In a fixed-update intervention, client updates and server state remain identical while the adversary changes only its report. The resulting weight change measures the direct sensitivity of the rule to unverified telemetry. The audit identifies a manipulable input and measures how much weight each rule concedes through it; rules that read no report act as controls whose invariance confirms that the intervention isolates the report channel.

We then introduce TrustFedGNN and its server rule, FU-Alignment. The server holds a small confidential audit graph, computes task and fairness reference gradients on it, clips client updates, and weights them by signed, temporally smoothed projections onto the combined reference. The aggregation path consumes no client-reported loss or disparity. Server anchoring itself is established by FLTrust~\citep{cao2021fltrust}; our contribution is the score dynamics (clipping before projection scoring and signed smoothing before rectification), whose effect we isolate by ablation, and an audit design that tests report dependence and adaptive-access boundaries.

The evidence supports three conclusions:
\begin{enumerate}[leftmargin=*]
    \item \textbf{Report-surface audit.} In a fixed-update audit, changing only the adversary's report changes its mean weight under six report-dependent rules, while CGSV, FLTrust, root-update FLTrust, and FU-Alignment produce bit-identical aggregates under a report falsified in every field except the sample count. An end-to-end run confirms the invariance for root-update FLTrust.
    \item \textbf{Controlled mechanism evidence.} Under a controlled fairness-poisoning protocol, FU-Alignment assigns less adversarial weight than root-update FLTrust on German (0.0630 versus 0.1892) and Bail (0.0451 versus 0.1959), with higher AUC on both datasets in the same matched runs. A factorial ablation attributes the reduction to clipping and signed smoothing; the fairness term has no detectable effect at the weight used.
    \item \textbf{Explicit boundaries.} An adaptive-access study shows that performance depends on dataset and adversarial access: under reference evasion FU-Alignment gives the adversaries the lowest weight and influence of the four rules compared and keeps the highest AUC on both datasets, and on Bail it has the lowest processed influence under sustained and intermittent poisoning; Section~\ref{subsec:rq4_boundary} identifies the access pattern under which this advantage does not hold. A known-label control supports AUC non-inferiority to FedAvg on Pokec-z observed labels, but does not support fairness superiority.
\end{enumerate}

\paragraph{Relation to prior work}
FSER and FTGD adapt ideas from our earlier FedFairGNN study~\citep{dang2026fedfairgnn} and are held fixed across relevant comparisons. They are not contributions of this paper. The present contribution lies in the trust-surface audit, the FU-Alignment server rule, and the evidence-backed characterization of its operating boundary.

Section~\ref{sec:related_work} positions the work against fairness-aware aggregation and server-anchored trust. Section~\ref{sec:methodology} defines the threat model and method. Section~\ref{sec:experiments} presents the evaluation, Section~\ref{sec:discussion} discusses deployment and limitations, and Section~\ref{sec:conclusion} summarizes the supported conclusions.
